\documentclass[border=2pt]{standalone}
\usepackage{amsmath, amsthm, amsfonts}
\usepackage{tikz-cd, kotex}
\usepackage{quiver}
\usepackage{Operators}
\usepackage{palette}
\usepackage{diagram-style}
\begin{document}

%%FIG 싸인 함수와 그 테일러 다항식 근사 (Taylor_Theorem-1.svg)
\begin{tikzpicture}
\tikzset{
  axes/.style={->, black!50, line width=0.4pt},
  sine/.style={accent1, line width=1.0pt},
  p1/.style={accent2, line width=0.9pt},
  p3/.style={black!50, line width=0.9pt},
  p5/.style={accent3, line width=0.9pt}
}
\draw[axes] (-3.8,0) -- (3.8,0);
\draw[axes] (0,-1.6) -- (0,1.6);
\begin{scope}
\clip (-3.15,-1.55) rectangle (3.15,1.55);
\draw[p1, domain=-3.15:3.15, samples=100] plot (\x, {\x});
\draw[p3, domain=-3.15:3.15, samples=100] plot (\x, {\x - \x*\x*\x/6});
\draw[p5, domain=-3.15:3.15, samples=100] plot (\x, {\x - \x*\x*\x/6 + \x*\x*\x*\x*\x/120});
\draw[sine, domain=-3.15:3.15, samples=200, smooth] plot (\x, {sin(\x r)});
\end{scope}
\end{tikzpicture}

\end{document}
